\documentclass{article}
\usepackage[utf8]{inputenc}
\usepackage[T1]{fontenc}
\usepackage[margin=2cm]{geometry}
\usepackage{hyperref}
\sloppy
\begin{document}
\section*{Historical source}
Encyclopaedia Britannica scan, source row 2464. Parse 3.5 extracted the text. Super 3.5 VL corrected it from an enlarged crop of each block. It can still contain recognition errors.
Page headers, footers, and page numbers are removed.

\subsection*{Untitled continuation (E0220)}
\textit{Continues an entry from an earlier page.}

C A B ( 2 ) C A B

which compose it : this cabbala is divided into three kinds, the gematrie, the notariicon, and the temura or themurah. The first whereof consists in taking the letters of a Hebrew word for ciphers or arithmetical numbers, and explaining every word by the arithmetical value of the letters whereof it is composed. The second sort of cabbala, called notariicon, consists in taking every particular letter of a word for an entire diction; and the third, called themurah, i. e. change, consists in making different transpositions or changes of letters, placing one for the other, or one before the other.

Among the Christians likewise, a certain fort of magic is, by mistake, called cabbala, which consists in using improperly certain passages of scripture for magic operations, or in forming magic characters or figures with stars and talismans.

Some visionaries, among the Jews, believe, that Jesus Christ wrought his miracles by virtue of the mysteries of the cabbala.

\subsection*{CABBALISTS (E0221)}
CABBALISTS, the Jewish doctors who profess the fludy of the cabbala.

In the opinion of these men, there is not a word, letter, or accent in the law, without some mystery in it. The Jews are divided into two general sects; the karaites, who refuse to receive either tradition or the talmud, or any thing but the pure text of scripture; and the rabbinists, or talmudists, who, besides this, receive the traditions of the ancients, and follow the talmud.

The latter are again divided into two other feets; pure rabbinifts, who explain the scripture, in its natural sense, by grammar, history, and tradition; and cabbalifts, who, to discover hidden mystical fenfes, which they suppose God to have couched therein, make use of the cabbala, and the mystical methods above mentioned.

\subsection*{CABECA (E0222)}
CABECA, or CABESSE, a name given to the finest silks in the East Indies, as those from 15 to 20 per cent. inferior to them are called barina. The Indian workmen endeavour to pass them off one with the other; for which reason, the more experienced European merchants take care to open the bales, and to examine all the skeins one after another. The Dutch distinguish two sorts of caceas; namely, the moor cacea, and the common cacea. The former is sold at Amsterdam for about 21½ schellingen Flemish, and the other for about 18½.

\subsection*{CABENDA (E0223)}
CABENDA, a port-town of Congo, in Africa, and subject to the Portuguese: E. long. 12°, and S. lat. 4°.

\subsection*{CABIDOS (E0224)}
CABIDOS, or CAVIDOS, a long measure used at Goa, and in other places of the East Indies belonging to the Portuguese, to measure stuffs, linens, \&c. and equal to ½ of the Paris ell.

\subsection*{CABIN (E0225)}
CABIN, in the sea-language, a small room, or apartment, whereof there are a great many in several parts of a ship; particularly on the quarter-deck, and on each side of the forecastle, for the officers of the ship to lie in.

The great cabin is the chief of all, and that which properly belongs to the captain, or chief commander.

\subsection*{CABINET (E0226)}
CABINET, or Cabinet, the most retired place in the finest part of a building, set apart for writing, studying, or preserving any thing that is precious.

A complete apartment consists of a hall, anti-chamber, chamber, and cabinet, with a gallery on one side. Hence we say, a cabinet of paintings, curiosities, \&c. CABINET also denotes a piece of joiner's workmanship, being a kind of press or chest, with several doors and drawers.

There are common cabinets of oak or of chestnut, varnished cabinets of China and Japan, cabinets of inlaid work, and some of ebony, or the like scarce and precious woods.

Formerly the Dutch and German cabinets were much esteemed in France, but are now quite out of date, as well as the cabinets of ebony, which came from Venice.

\subsection*{CABIRI (E0227)}
CABIRI, a term in the theology of the ancient Pagans, signifying great and powerful gods; being a name given to the gods of Samothracia. They were also worshipped in other parts of Greece, as Lemnos and Thebes, where the cabiria were celebrated in honour of them: these gods are said to be, in number, four, viz. Axieros, Axioerfa, Axioerfus, and Cafinilus.

\subsection*{CABIRIA (E0228)}
CABIRIA, festivals in honour of the cabiri, celebrated in Thebes and Lemnos, but especially in Samothracia, an island consecrated to the cabiri. All who were initiated into the mysteries of these gods, were thought to be secured thereby from storms at sea, and all other dangers. The ceremony of initiation was performed by placing the candidate, crowned with olive-branches, and girded about the loins with a purple ribband, on a kind of throne, about which the priests, and persons before initiated, danced.

\subsection*{CABLAN (E0229)}
CABLAN, the name of a kingdom and city of India, beyond the Ganges.

\subsection*{CABLE (E0230)}
CABLE, a thick, large, strong rope, commonly of hemp, which serves to keep a ship at anchor.

There is no merchant-ship, however weak, but has at least three cables; namely, the chief cable, or cable of the fleet-anchor, a common cable, and a smaller one.

Cable is also said of ropes, which serve to raise heavy loads, by the help of cranes, pullies, and other engines. The name of cable is usually given to such as have, at least, three inches in diameter; those that are less are only called ropes of different names, according to their use.

Every cable, of what thickness soever it be, is composed of three strands; every strand of three ropes; and every rope of three twists: the twist is made of more or less threads, according as the cable is to be thicker or thinner.

In the manufacture of cables, after the ropes are made, they use flicks, which they pass first between the ropes of which they make the strands, and afterwards between the strands of which they make the cable, to the end that they may all twist the better, and be more regularly wound together; and also, to prevent them from twining or intangling, they hang


\section*{Source record}
\begin{description}
\item[Scan ID] \texttt{britannica:309f765b1ecc83e0a7f7762e44b4a927a7555443:train:2464}
\item[Image file] \texttt{page-2464-24f450326eee.jpg}
\item[Image SHA-256] \texttt{24f450326eee13b008306d2b355bd49b295db96603186268db20b85552cf958e}
\item[Dataset] NationalLibraryOfScotland/encyclopaedia\_britannica\_illustrated at revision \texttt{309f765b1ecc83e0a7f7762e44b4a927a7555443} (CC0-1.0)
\item[Extraction] model \texttt{nvidia/nemotron-3.5-parse}, raw reply \texttt{extract/raw/page-2464-24f450326eee.json}
\item[Raw reply SHA-256] \texttt{ca4563d6e384df719eaa01880d3377a6baf8c1433206faef512c45b6b77735c3}
\item[Correction] model \texttt{nvidia/Nemotron-3.5-Super-VL}, block crops, raw replies \texttt{correct/raw/crop\_nothink/page-2464-24f450326eee-b*.json}; super\_vl: 23
\item[Running header] none
\end{description}
\end{document}
